% chapters/ch02.tex —— 第 2 章
% 体例：具体失败 → 设计问题 → 机制解释 → 对照实验 → 条件变化 → 选择原则
% 数字身份纪律：本章实测数字全部来自演示运行 CH02DEMO-20260923T153152Z-56208b
% （原始记录 experiments/runs/CH02DEMO-20260923T153152Z-56208b/）；本章无外部引用数字。
\chapter{三层时间尺度下的自进化框架}
\label{ch02}

\chaptergoal{一次执行和一次修改，怎么接成一条能追责的过程？}
\chapterprops{P1、P1b（四格比较需要 $A_k$ 与可更新集合的精确定义）；P6、P7（退化与元层改进需要版本与血缘）。本章建立这些问题所需的系统对象。}
\chapterfig{图~\ref{fig:ch02-loop}：一次更新的闭环——四个角色各自负责什么，验证器 $\mathcal V$ 站在哪里}
\chapterpitfall{三次都是通过，但后两次打开文件时，答案已经在那里了}

% ---------- 1. 具体失败 ----------
\begin{failurebox}
同一道闭区间求和题，同一个 Agent 和同一份系统提示，在两种条件下各做三次。一组每次在新目录里放入题目给定的文件，另一组三次共用一个目录，只在第一次放入初始文件。两组验收都是 \measured{3/3}。

差别藏在 \texttt{main.py} 的一行代码里。初始程序用 \texttt{range(a, b)} 求和，漏掉了右端点。在新目录里，Agent 每次都读到这段程序，再把上界改成 \texttt{b + 1}。在共用目录里，只有第一次做了这项修改。第二、三次打开的已经是修好的程序，它们读完文件便运行验收，没有再写文件。

运行汇总也记下了每次开始时清单内的文件文本是否与题目给定的文本一致。新目录组为 \measured{3/3}，共用目录组为 \measured{1/3}。后两次继承的是工作目录里的答案。六份记录足以看出解题过程的差异，至于这种残留会怎样影响其他任务的通过率，还需要在那些任务上比较。
\end{failurebox}

\designquestion{系统的“下一次”从哪里开始？这句话没定义清楚之前，任何“变好了”都无法归因。}

第 1 章把“留下什么”变成了一个可检查的问题，但还没有给出检查它的对象。
本章补上这套对象：三层时间尺度各自由什么变量承载，状态分成哪三类，
一次更新由哪四个角色完成，候选版本经过什么决策才成为新版本，以及版本之间靠什么串成血缘。
这些定义为 E01 四格提供共同的描述语言，但同一套 $A_k$ 定义与相同 episode 配额还不构成因果对照。历史四格在执行循环、经验来源与总成本上存在差异，第3章据此限定其结果解释。

% ---------- 2. 机制解释 ----------
\section{任务内动作、单任务轨迹与版本更新}
\label{sec:ch02-layers}

一个自进化系统在三个时间尺度上同时活动。第 1 章给出过三层的名字，本节给出它们各自承载的变量与判据。

\begin{table}[htbp]
\centering
\caption{三层时间尺度：变了什么、凭什么说它变了、最容易和什么混淆}
\label{tab:layers}
\begin{tabular}{>{\raggedright\arraybackslash}p{1.6cm}>{\raggedright\arraybackslash}p{2.0cm}>{\raggedright\arraybackslash}p{3.6cm}>{\raggedright\arraybackslash}p{3.4cm}>{\raggedright\arraybackslash}p{3.0cm}}
\toprule
层 & 变量 & 变化的东西 & 判定它真的变了的检查 & 典型混淆 \\
\midrule
步 $t$ & $H_{k,e,t}$ & 工作上下文：消息、工具结果、中间产物 & 清空上下文后该变化是否消失 & 上下文变长被当成持久更新 \\
回合 $e$ & $\tau_{k,e}$ & 一次任务尝试的轨迹，可含反思与重试 & 哪些步骤属于这次尝试，哪些临时状态随回合复位 & 把归档的轨迹当成已被下一回合读取的记忆 \\
版本 $k$ & $A_k$ & 系统的持久组件，评价设置另列 & 核对持久组件的变更与后续加载记录 & 环境残留被当成持久更新 \\
\bottomrule
\end{tabular}
\end{table}

前两条检查用于定位变化落在哪一层，第三条用于判断能否进一步评价变化的效果：

\begin{enumerate}
  \item \textbf{留存判据}：任务结束后，这个变化还在不在？不在，它最多属于 $e$ 层。
  \item \textbf{起点判据}：下一个任务的出发点有没有因为它而不同？没有，就说明它没有改变这次任务的起点；是否影响更晚的任务，还需检查后续加载与使用记录。
  \item \textbf{独立检验判据}：这个变化能不能被一个不参与修改的机制检验？不能，就尚无独立证据支持它“变好了”；这不否定更新本身可能已经发生。
\end{enumerate}

第三条需要另外准备评价证据，不能由前两条替代。本章开头的演示里，条件二的后两次运行满足留存判据（文件确实留着）
与起点判据（下一次确实从改过的文件出发），但没有证明受控持久组件带来了收益。
这次历史演示的归档编号为 \path{CH02DEMO-20260923T153152Z-56208b}。汇总保存了清单内文件的文本比对结果，轨迹则记下了 Agent 实际读写的内容。两者都能帮助解释这道题的起点，但无法据此还原整个目录的原始字节。

在实验台的 \path{experiments/lab} 目录里，可以沿一条回合记录检查登记过的写入。\path{trace.py} 定义的 \texttt{EpisodeTrace} 把执行步骤存入轨迹，并用 \path{persistent_writes} 字段登记持久写入。静态基线在 \path{static_agent.py} 中检查这个字段，只要有一笔写入就报错。即使字段为空，仍须另查工作目录、缓存与外部服务是否留下了影响下一回合的状态。

\section{持久状态、工作上下文与环境状态}
\label{sec:ch02-states}

“状态”在 Agent 的语境里至少指三种不同的东西，把它们混在一起是本章开头那个假象的根源。

\begin{table}[htbp]
\centering
\caption{三类状态：归谁所有、何时消失、越界了会怎样}
\label{tab:states}
\begin{tabular}{>{\raggedright\arraybackslash}p{2.6cm}>{\raggedright\arraybackslash}p{3.2cm}>{\raggedright\arraybackslash}p{3.4cm}>{\raggedright\arraybackslash}p{4.4cm}}
\toprule
类别 & 归属 & 消失时机 & 越界的后果 \\
\midrule
工作上下文 $H_{k,e,t}$ & 本回合可见的消息与工具结果 & 随回合复位，轨迹可另行归档 & 把本次提示中已有的信息当成下一次也能调用的能力 \\
持久组件 $\theta,P,M,L,G$ & 系统版本 $A_k$ & 按受控更新或回退改变 & 改动未经 Gate 就生效，收益与损失都无法归因 \\
环境状态 & 文件系统、数据库、缓存、外部服务 & 依协议决定复位、固定或记录 & \textbf{改变下一次运行的起点，却不出现在任何版本记录里} \\
\bottomrule
\end{tabular}
\end{table}

共用目录里留下的答案对应表~\ref{tab:states} 的“环境状态”一行。它既不属于 $A_k$，也不属于下一回合刚建立的工作上下文，
它只是“恰好还在那里”。它改变下一次运行的起点，可能影响通过率，仅凭分数曲线难以把它与受控持久更新区分，
在这个演示中，它没有经过持久更新的发布过程。其他系统可以把环境快照纳入版本控制；此时应明确扩大后的系统边界。

\begin{sloppypar}
历史演示改变的是\textbf{工作目录的复位方式}。新目录可以避免读到上一次的答案，却管不到候选进程能用哪些宿主资源。当前实验台的文件工具允许读取验收断言，会拒绝直接写入受保护的验收文件。执行候选还需要操作系统级隔离；目前缺少这一执行器，演示入口在创建模型后端之前就会停止。
\end{sloppypar}

由此得到一条可写进协议的复位判据：
\textbf{每次运行前按协议复位任务目录，核对实际文件的相对路径和内容。}若协议要求逐字节相同，还须检查文件字节，并识别清单外多出的文件。
比对失败说明起点与协议预期不符。需要先确认差异来源；在排除它对结果的影响之前，不能把全部分数差归给受控更新。
公共协议第8条规定了参考答案与验收器的隔离目标。本章的例子说明，即使验收器始终相同，仍须检查求解从哪个\textbf{起点}出发。

% figures/ch02-loop.tex —— 第 2 章核心图（全书引用入口）
% 用法：在 book/chapters/ch02.tex 中 % figures/ch02-loop.tex —— 第 2 章核心图（全书引用入口）
% 用法：在 book/chapters/ch02.tex 中 % figures/ch02-loop.tex —— 第 2 章核心图（全书引用入口）
% 用法：在 book/chapters/ch02.tex 中 \input{figures/ch02-loop}
% 图主体在 ch02-loop-body.tex，便于单独编译与复用。
\begin{figure}[htbp]
\centering
\input{figures/ch02-loop-body}
\caption{一次更新的闭环：执行器产生轨迹，反馈器产生信号，修改器产生候选版本，
Gate 决定候选暂不发布、影子运行还是灰度发布。$\mathcal V$ 默认冻结，按固定的任务划分与验收规则，为执行和候选比较提供参照。}
\label{fig:ch02-loop}
\end{figure}

% 图主体在 ch02-loop-body.tex，便于单独编译与复用。
\begin{figure}[htbp]
\centering
% Vertical layout preserves final-size labels.
\begin{tikzpicture}[
 box/.style={draw,rounded corners=2pt,minimum height=11mm,text width=49mm,align=center,font=\small},
 small/.style={draw,rounded corners=2pt,text width=27mm,minimum height=12mm,align=center,font=\small},
 arrow/.style={-{Stealth[length=1.8mm]},thick}
]
\node[box,fill=blue!7] (ak) at (0,0) {$A_k$：已发布版本};
\node[box,fill=gray!7] (run) at (0,-1.7) {执行器\\$\tau_{k,e}=\mathrm{Run}(A_k,T)$};
\node[box,fill=gray!7] (fb) at (0,-3.4) {反馈器：产生信号$f$};
\node[box,fill=orange!9] (upd) at (0,-5.1) {修改器：产生候选\\$\widetilde A_{k+1}=\mathrm{Update}(A_k,\tau,f)$};
\node[box,fill=green!8] (gate) at (0,-6.8) {Gate：三级发布决策};
\draw[arrow] (ak)--(run);\draw[arrow] (run)--(fb);
\draw[arrow] (fb)--(upd);\draw[arrow] (upd)--(gate);
\node[box,fill=red!5,draw=red!55!black] (v) at (7,-2.2) {$\mathcal V$默认冻结\\任务划分、验收器、最终测试集\\不在Updatable内（公理1）};
\draw[arrow,red!55!black] (v.west)--++(-.8,0)|-(run.east);
\draw[arrow,red!55!black] (v.south)--(7,-4.2)--(3.8,-4.2)--(3.8,-6.8)--(gate.east);
\node[small,fill=red!5] (reject) at (-3.4,-8.7) {reject\\保留候选档案};
\node[small,fill=yellow!10] (shadow) at (0,-8.7) {shadow\\影子运行};
\node[small,fill=green!8] (accept) at (3.4,-8.7) {accept($r$)\\按比例灰度发布};
\draw[arrow] (gate.south west)--(reject.north);
\draw[arrow] (gate)--(shadow);\draw[arrow] (gate.south east)--(accept.north);
\node[box,fill=blue!7] (next) at (7,-6.0) {$A_{k+1}$：被接受的下一版};
\draw[arrow,dashed] (accept.east)-|(next.south);
\end{tikzpicture}

\caption{一次更新的闭环：执行器产生轨迹，反馈器产生信号，修改器产生候选版本，
Gate 决定候选暂不发布、影子运行还是灰度发布。$\mathcal V$ 默认冻结，按固定的任务划分与验收规则，为执行和候选比较提供参照。}
\label{fig:ch02-loop}
\end{figure}

% 图主体在 ch02-loop-body.tex，便于单独编译与复用。
\begin{figure}[htbp]
\centering
% Vertical layout preserves final-size labels.
\begin{tikzpicture}[
 box/.style={draw,rounded corners=2pt,minimum height=11mm,text width=49mm,align=center,font=\small},
 small/.style={draw,rounded corners=2pt,text width=27mm,minimum height=12mm,align=center,font=\small},
 arrow/.style={-{Stealth[length=1.8mm]},thick}
]
\node[box,fill=blue!7] (ak) at (0,0) {$A_k$：已发布版本};
\node[box,fill=gray!7] (run) at (0,-1.7) {执行器\\$\tau_{k,e}=\mathrm{Run}(A_k,T)$};
\node[box,fill=gray!7] (fb) at (0,-3.4) {反馈器：产生信号$f$};
\node[box,fill=orange!9] (upd) at (0,-5.1) {修改器：产生候选\\$\widetilde A_{k+1}=\mathrm{Update}(A_k,\tau,f)$};
\node[box,fill=green!8] (gate) at (0,-6.8) {Gate：三级发布决策};
\draw[arrow] (ak)--(run);\draw[arrow] (run)--(fb);
\draw[arrow] (fb)--(upd);\draw[arrow] (upd)--(gate);
\node[box,fill=red!5,draw=red!55!black] (v) at (7,-2.2) {$\mathcal V$默认冻结\\任务划分、验收器、最终测试集\\不在Updatable内（公理1）};
\draw[arrow,red!55!black] (v.west)--++(-.8,0)|-(run.east);
\draw[arrow,red!55!black] (v.south)--(7,-4.2)--(3.8,-4.2)--(3.8,-6.8)--(gate.east);
\node[small,fill=red!5] (reject) at (-3.4,-8.7) {reject\\保留候选档案};
\node[small,fill=yellow!10] (shadow) at (0,-8.7) {shadow\\影子运行};
\node[small,fill=green!8] (accept) at (3.4,-8.7) {accept($r$)\\按比例灰度发布};
\draw[arrow] (gate.south west)--(reject.north);
\draw[arrow] (gate)--(shadow);\draw[arrow] (gate.south east)--(accept.north);
\node[box,fill=blue!7] (next) at (7,-6.0) {$A_{k+1}$：被接受的下一版};
\draw[arrow,dashed] (accept.east)-|(next.south);
\end{tikzpicture}

\caption{一次更新的闭环：执行器产生轨迹，反馈器产生信号，修改器产生候选版本，
Gate 决定候选暂不发布、影子运行还是灰度发布。$\mathcal V$ 默认冻结，按固定的任务划分与验收规则，为执行和候选比较提供参照。}
\label{fig:ch02-loop}
\end{figure}


\section{三个索引不能排成同一根时钟}
\label{sec:ch02-index-boundaries}

步、回合与版本是三种记录单位，并不意味着每结束一个回合就一定产生一个新版本。一个系统可以连续完成许多任务而不更新，也可以在一次任务中提出多个候选，最后全部拒绝。若机械地把任务编号当成版本编号，记录会制造并不存在的更新；若只在成功更新时给任务编号，被拒绝候选所依赖的任务又可能消失。两种编号都应保留，它们之间通过事件关联。

\subsection{执行步描述发生顺序，版本描述生效条件}

在一个固定版本的回合内，步索引 $t$ 可以按真实发生的事件前进：模型返回文本，工具接受调用，工具返回结果，执行器决定继续或结束。模型文本提出一个动作，并不等于工具已经执行该动作。比如回复包含格式错误的参数，执行器可能在调用工具之前就拒绝它。此时轨迹里应该有提议与拒绝，不能凭模型文本生成一条虚假的工具成功记录。

把这两层分开，还能解释为什么模型调用次数不等于工具调用次数。一次模型回复可能没有合法动作；一个合法动作可能只读取文件；协议也可能允许工具内部完成多项工作。研究若以调用次数为预算，需要明确计数的是模型请求、工具入口还是更底层的外部请求。它们都可以形成事件序列，却没有天然的一一对应关系。

版本索引 $k$ 则回答某项更新何时获得生效资格。候选产生并不使 $k$ 自动前进，验证失败也不应改变当前发布对象。即使选择器接受了候选，已经开始的回合是否立即切换，还要看协议。固定回合版本的设计让整条轨迹面对同一组组件；允许热切换的设计更灵活，但需要按步记录组件切换，并另外解释切换前后信息如何传递。

为描述后一种情况，可以给每个动作附上实际加载的版本标识 $v(e,t)$。如果协议规定回合内固定版本，那么对同一个 $e$，$v(e,t)$ 在所有执行步保持一致；若它发生变化，记录就暴露了协议违例或一个需要声明的切换。这里的标识用于追踪来源，不要求每个执行步都复制全部模型与技能文件。

固定版本也不等于固定输出。同一组组件面对不同上下文会产生不同动作，随机采样也可能产生不同结果。因此“动作不同”不能独立证明版本发生变化；“文件相同”也不能保证执行条件相同，因为远程服务、工具环境与输入仍可能不同。版本清单只是复现条件的一部分，需要与任务及环境记录一起阅读。

\subsection{回合结束，是协议边界而不是自然现象}

编程题容易定义回合：从加载题目开始，到提交答案、超出预算或装置失败结束。持续运行的助手则未必有同样清楚的边界。一段对话结束以后，稍后的补充指令究竟属于原任务，还是一个新任务，取决于系统如何划分。划分不清，就无法判断一次总结是任务内压缩，还是已经进入跨任务记忆。

因此协议应说明终止事件和继承规则。超时结束的回合是否允许产生经验；失败任务的临时文件是否可以进入候选；用户撤销任务后是否还继续后台整理；这些都影响下一任务能接触的信息。答案不必一律为否，但必须在比较之前确定，不能看见结果以后才选择一种有利的边界。

回合结束还不代表所有外部效果已经结束。一个异步工具请求可能仍在运行，后台索引可能稍后才刷新。如果下一回合在这些效果落定之前开始，就需要明确哪些状态被继承，哪些请求仍待确认。只有“任务完成”的文字消息，不能证明外部过程已经静止。对无法停止的环境，协议应记录可观察状态与时间窗口，而不是声称已经彻底复位。

这也解释了为什么本章开头的新建目录检查只适用于它所覆盖的边界。它能够发现任务文件残留，却不能证明远程服务没有缓存，也不能证明共享账户的额度没有变化。一个有限检查有自己的价值；把它扩大成整个环境完全一致的证明，反而会让读者误用结果。

\subsection{反馈可以迟到，更新不能因此失去来处}

有些反馈在任务结束时立即得到，例如本地验收退出码；另一些反馈要等外部结果发生后才能得到。迟到反馈到达时，系统可能已经发布了多个新版本。把它直接挂到“当前版本”名下，会使一条旧任务的结果看起来像是在评价最新版本。正确的关联首先应指向产生该结果的回合及其实际执行条件。

迟到反馈仍可以用于产生新的候选，但应同时记录两件事：证据来自哪个历史版本，新候选准备修改哪个当前父版本。二者可能不同。历史问题也许已经被中间某次更新解决，也可能在新版本上以另一种方式出现；修改器需要先检查当前对象，不能把针对旧实现的修补直接套上去。

这个问题不只出现在大型线上系统。一个实验在中断后恢复，先读到之前尚未汇总的验收结果，也会面临同样的关联要求。用文件修改时间猜测属于哪次尝试很脆弱；显式回合标识和请求标识更容易复核。如果缺失的关联已经无法恢复，应把该记录标为归属不明，避免把它补进最方便的一组数据。

\section{用复位矩阵声明系统边界}
\label{sec:ch02-reset-matrix}

“每次清空环境”听起来明确，实际可能清得太少，也可能清得太多。清得太少，旧答案进入新任务；清得太多，原本要研究的持久记忆也被删掉，实验退化成没有学习的系统。需要复位的不是抽象的“所有状态”，而是协议规定不应跨越当前边界的状态。

\subsection{先列对象，再给每个对象一种处理规则}

对本书的编程题，可以区分题目初始文件、任务内生成文件、发布技能、候选技能、检索缓存与外部请求记录。任务初始文件按新题重建；任务内文件在保存必要轨迹后清理；发布技能按选定版本加载；候选技能保留在未发布区域；缓存根据它是否属于被研究系统来决定复位还是纳入版本；外部请求记录则需要保留，以便避免丢失成本与效果证据。

这组规则构成一张复位矩阵：每类对象在任务开始、任务结束、候选接受与恢复时分别如何处理。矩阵的目的不是增加格式，而是防止同一对象在不同模块中被赋予冲突含义。运行器认为缓存是临时状态，检索器却依赖缓存跨任务积累信息，两者即使各自工作正常，也会让实验的系统边界失去一致性。

\begin{table}[htbp]
\centering\small
\caption{编程题设计示例的复位矩阵；这是建议的对象规则，不是新增实测}
\label{tab:ch02-reset-matrix}
\begin{tabular}{>{\raggedright\arraybackslash}p{2.0cm}>{\raggedright\arraybackslash}p{4.0cm}>{\raggedright\arraybackslash}p{3.8cm}>{\raggedright\arraybackslash}p{4.0cm}}
\toprule
对象 & 新任务开始 & 任务结束 & 接受候选或恢复 \\
\midrule
题目与任务产物 & 按该题初始清单重建 & 先留轨迹，再清理临时物 & 不继承上一题答案 \\
已发布技能 & 加载选定版本 & 按版本保留 & 切换后确认实际加载对象 \\
未发布候选 & 默认不进入任务输入 & 连同拒绝原因留档 & 只发布已经绑定证据的工件 \\
检索缓存 & 按声明策略固定或清空 & 记录保留内容与作用范围 & 防止旧缓存遮蔽新版本 \\
外部请求记录 & 读取未确认请求状态 & 保留成本与效果记录 & 不因本地恢复而删去历史消耗 \\
\bottomrule
\end{tabular}
\end{table}

这张表没有把所有外部环境都纳入受控组件。若某个服务不能复位，就把它列为受观察但不能固定的条件，并说明它可能怎样影响比较。用明确限制替代虚假的完全控制，通常能使结果的适用范围更清楚。研究者随后可以决定是否增加一个离线替代环境，但替代环境的结果也需要单独标识。

\subsection{检查起点时，要比对期望状态而不是旧快照}

上一回合的起点不一定是下一回合的正确起点，因为题目本来就可能不同。正确检查对象应是“新任务定义加上选定版本”共同规定的期望状态。题目文件应该匹配新题；发布技能应该匹配本次选择的版本；不允许继承的目录应该为空。若把所有文件都与上一题的快照逐字节比对，合法的任务变化反而会被判成污染。

期望状态也可以允许某些差异。例如运行编号、临时路径和时间戳天然会变化。忽略这些字段需要明确规则，不能在看到比对失败以后临时把差异加入例外。否则，检查器可能逐渐学会忽略所有重要变化。例外规则本身也是实验配置的一部分，修改后应重新说明哪些历史检查仍然适用。

文件清单与内容校验各解决一类问题。只比对已有文件内容，可能漏掉多出来的答案文件；只比对文件名，又可能漏掉同名文件被改写。因此检查应同时考虑期望对象、实际对象与内容差异。对于数据库或索引，不必强求用文件方式检查，但应提供同等明确的逻辑状态摘要或版本入口，并声明摘要没有覆盖的内容。

\subsection{恢复之后，重新建立的是哪条执行链}

恢复有时意味着重新执行尚未完成的任务，有时意味着只恢复持久组件供下一任务使用。前者需要处理已经发生的工具效果，后者主要需要确认组件与依赖能加载。两种恢复若共用一个“resume”标志，却没有不同的说明，容易让使用者以为任务会从中断的精确位置无缝继续。

设一次工具请求已发出，但确认消息丢失。重新执行请求可能完成原本未完成的动作，也可能重复已经发生的动作。这个不确定性不能靠清空本地状态消除。对有外部写入的工具，需要使用服务提供的查询或去重能力确认状态；若没有这种能力，就应显式保留未知效果，避免把重新执行包装成没有副作用的恢复。

纯本地、可重复的检查通常更容易恢复，但也必须确认输入未变。对同一个固定工件重新运行验收，可以得到一条新的检查记录；它不能改写中断前那次检查的完成状态。把两次尝试分开保存，读者才能知道哪次耗时和成本已经发生，哪次真正产生了可用结果。

恢复到旧组件后，如果继续沿用新组件产生的缓存或经验，也可能得到一个未被评价的组合。完整恢复未必要求删除所有新信息，但必须明确哪些信息继续保留，以及这是否仍是原本声明的旧版本。如果保留改变了后续行为，最好把它视为新的组合状态进行检查，不能仅凭组件目录名称恢复为旧值，就宣称整个系统已经回到过去。

\subsection{从矩阵发现设计遗漏}

填写矩阵时，如果某个对象无法确定归属，往往说明系统设计还藏着一个未回答的问题。检索缓存是否属于长期记忆，取决于它会不会影响新任务、是否被允许跨任务保留，以及是否接受受控更新。对象叫作“缓存”并不能豁免版本与成本记录；对象叫作“记忆”也不能保证它真的进入了后续执行。

矩阵还可以暴露权限矛盾。某个验收文件被声明为不可修改，执行器却可以通过更通用的文件工具写入它，那么上层约定尚未被底层权限落实。保护其中一个入口不足以证明其他入口不存在。对本书当前原型，工具层保护与操作系统级隔离应分别描述，后者没有实现就不应借用前者的名称。

最后，边界需要与研究问题一起变化。研究跨任务记忆时，应让规定的记忆继续保留；研究无持久物的静态基线时，应阻止相同信息从其他路径进入。两组环境不必在所有字节上完全相同，因为被研究的持久组件本来就是处理差异；除被研究的处理差异之外，需要保持一致的是任务生成规则、其余允许的信息来源和预先声明的比较条件。这个区别是从“每次清空”走向可解释实验的关键。


复位检查本身也应接受有针对性的验证。可以在离线副本中故意放入一个上一任务的答案文件，检查它是否被发现；再改变发布技能的内容而保留原文件名，确认内容检查不会只看名称；最后保留一项协议允许的记忆，确认复位过程不会误删它。这些是检验检查器的建议案例，本章并未把它们报告成已经运行的新实验。

只验证“干净目录能通过”，不足以说明检查器能拒绝污染目录；只验证“污染目录被拒绝”，也不足以说明合法持久物能够保留。两类检查对应漏检与误拒，应该分别观察。发现检查器错误后，受影响的是哪些运行、是否需要重算起点一致性，都应回到已保存的对象清单中判断。若当时没有保存足够材料，就只能承认相关历史条件无法复核，不能靠今天的检查器为昨天补发通过记录。

\section{执行器、反馈器、修改器与选择器}
\label{sec:ch02-roles}

一次更新在图~\ref{fig:ch02-loop} 里被拆成四个角色，它们承担不同职责，具体实现可以共享程序或模型。

\begin{definition}[一次更新的四个角色]
执行器在给定任务上运行 $A_k$ 并产生回合记录 $\tau_{k,e}$；
反馈器产生信号 $f$；
修改器由搜索空间 $\mathcal S$（改哪个组件）与优化算法 $H$（怎么选）组成，
按 $\operatorname{Update}(A_k,\tau,f)$ 产生候选版本 $\widetilde A_{k+1}$；
选择器（Gate）决定候选是否成为 $A_{k+1}$。
\end{definition}

\begin{equation}
\tau_{k,e}=\operatorname{Run}(A_k,T_{k,e}),\qquad
\widetilde A_{k+1}=\operatorname{Update}(A_k,\tau_{k,e},f)
\label{eq:update}
\end{equation}

修改器的 $(\mathcal S,H)$ 二分不是本书提出的，它来自反馈环综述对优化器的分解\cite{self-evolving-comprehensive-2508}。
本书沿用这个分解，并增加一条约束：\textbf{搜索空间中允许改变的配置，只能涉及可更新集合 Updatable 中的组件}。
可更新集合是允许被修改的组件集合，$\mathrm{Updatable}\subseteq\{\theta,P,M,L,G\}$。
收紧之后，“改了什么”就有了一个可以指的位置，而不是一句“系统更新了”。

表~\ref{tab:roles} 把四个角色与本书实验台的落点对齐。落点这一列不是为了展示代码，
而是为了明确哪些职责应当能够分别检查。E01 四格在修改范围上有如下差别：
格 1 不做持久更新，格 3 的注意事项只在本任务内有效，格 4 把经验写进 $M$ 并跨任务保留。历史实现还改变了执行过程，因此不能把它描述成只替换修改器的单因素实验。

\begin{table}[htbp]
\centering
\caption{四个角色：输入、输出、在本书实验台上的落点、坏掉时的症状}
\label{tab:roles}
\begin{tabular}{>{\raggedright\arraybackslash}p{1.5cm}>{\raggedright\arraybackslash}p{2.4cm}>{\raggedright\arraybackslash}p{2.4cm}>{\raggedright\arraybackslash}p{3.6cm}>{\raggedright\arraybackslash}p{3.1cm}}
\toprule
角色 & 输入 & 输出 & 实验台落点 & 坏掉时的症状 \\
\midrule
执行器 & $A_k$、任务 $T$ & $\tau_{k,e}$ & \texttt{static\_agent.py} 的 \texttt{run\_episode} & 轨迹不完整，后面全都无从诊断 \\
反馈器 & 执行结果 & 信号 $f$ & 旧版 \texttt{task.verify} 执行验收命令 & 信号与真实目标不一致，改进方向被带偏 \\
修改器 & $A_k,\tau,f$ & $\widetilde A_{k+1}$ & \texttt{e01\_arms.py} 的写入与适配 & 改对了分数，改错了组件 \\
选择器 & $\widetilde A_{k+1}$、验证结果 & 部署决定 $\mathcal D_{k+1}$ & 各实验分别核对；历史E01没有经验质量准入 & 未经检查的候选生效，拒绝理由缺失 \\
\bottomrule
\end{tabular}
\end{table}

反馈器与验证器 $\mathcal V$ 不是一回事，但两者常被混用。
$\mathcal V$ 是任务划分、验收器与最终测试集的合称，它\textbf{默认冻结}，不在 Updatable 内（公理 1）；
反馈器只是产生信号 $f$ 的部件，它的信号可以来自 $\mathcal V$，也可以来自模型自评、人类标注或其他系统。
把两者分开之后，“反馈质量”与“验证是否被污染”才是两个可以分别讨论的问题——第 4 章处理前者，第 14 章处理后者。

公理 1 的价值在这里第一次显出结构性的后果。收益估计写成

\begin{equation}
\widehat\Delta_k(b)=\frac{1}{n}\sum_{i=1}^{n}\bigl[q(A_{k+1},x_i;b)-q(A_k,x_i;b)\bigr]
\label{eq:delta}
\end{equation}

时，任务集 $\{x_i\}$ 与验收器 $q$ 都是$\mathcal V$ 的一部分。
如果比较前后同时更换 $\mathcal V$，任务与评分函数也可能改变：
系统既能改自己的答案，也能改判分的方式与判分的题目，式~\eqref{eq:delta} 不再比较同一个东西。
\textbf{差值仍可以计算，但不再单独估计系统版本改变对同一任务、同一评分目标的作用。}
因此，若把 $\mathcal V$ 纳入 Updatable，就需要单独设置干预组，并保留共同的外部评价条件；第14章讨论这种设计与验收污染的区别。第12章主要研究修改器本身的更新，它不自动意味着验证器也被修改。

公理 2 在这一节的落点是“差异定位”。
可更新集合每扩大一档，就多出一类“改动可能藏身的地方”：
若其余条件保持固定，可更新集合只有 $M$ 时，声明的持久更新对象就是记忆；把 $L$ 放进来，技能库的行为变更也要纳入差异定位。
历史 E01 的格 4 只把记忆 $M$ 列为显式更新对象。训练时，\texttt{extract\_lesson} 根据最终代码与验收结果生成经验文本；训练循环再把文本和来源标识追加到记忆列表。要排除未登记的变化，还须另查依赖与环境。
扩张的代价因此变得具体：把 $L$ 或 $G$ 加进可更新集合，就要新造一条写入路径，
并为这条路径重新回答一遍“凭什么认为它有用”。

\section{三级 Gate：reject、shadow 与 accept}
\label{sec:ch02-gate}

候选版本产生之后，还差一个决定。二分（接受 / 拒绝）不够用，原因很具体：
有一大类候选值得观察但不值得发布——它通过了验证，可谁也不知道它在真实流量上表现如何。

\begin{equation}
\mathcal D_{k+1}=\begin{cases}
A_k, & \text{reject：保留诊断与候选记录}\\[2pt]
A_k\ \text{发布},\ \widetilde A_{k+1}\ \text{影子运行}, & \text{shadow}\\[2pt]
\text{按灰度比例 } r\in(0,1]\ \text{混跑}, & \text{accept}(r)
\end{cases}
\label{eq:gate}
\end{equation}

这里 $\mathcal D$ 表示部署决定，不是式~\eqref{eq:ak} 的组件元组；观察或混跑需要同时记录新旧版本和路由。三级还改变了候选获得的执行权限。reject 保留候选及诊断；
shadow 让候选在线上并行跑，它产生的结果不计入发布版本，在隔离其外部写入的前提下，可用于观察真实分布上的表现；
accept($r$) 才让候选承担流量，$r$ 是灰度比例。
“影子更新 + 定期合入”这种工程上常见的做法，在二分 Gate 里没有位置，在三级 Gate 里是 shadow 到 accept 的一段。

第 13 章的触发条件说的正是这段路：触发判据决定的是\textbf{什么时候把 shadow 升成 accept}，
这是触发的一种用途；产生候选、启动验证与停止发布也各有触发条件。需要明确条件控制的是哪个阶段，避免把“已经在影子运行”和“还没开始”混成一件事。

\begin{sloppypar}
三级 Gate 是本书用于分析发布决定的框架，不能整体当成历史实验已经实现的功能。
历史 E01 的格4每轮根据最终代码与任务验收结果生成经验，随后直接追加到经验列表；任务失败并不阻止写入。
E05 的技能准入另有协议，不能把某个实验的拒绝条件套到全书。
shadow 与灰度发布需要流量路由、效果隔离与恢复设施，第13、14章讨论这些要求；讨论不等于已在实验台实现。
没有 shadow 只说明没有这一观察阶段，不能据此推断发布只能发生一次，或系统一定不能回滚。
\end{sloppypar}

\subsection{任务验收不是经验准入}
\label{sec:ch02-source-verdict-admission}

这个区别可以直接沿历史 E01 的实现检查。以下两个文件位于实验台的 \texttt{lab} 目录。

\begin{sloppypar}
\texttt{e01\_arms.py} 中的经验提取函数，把任务陈述、最终代码和“通过／未通过”交给模型，返回经过长度截断的文本；
\texttt{run\_e01.py} 的训练阶段随后把文本加入列表，并保存来源判定字段。
这里的判定字段描述经验来自一次怎样的任务运行，代码没有再对经验陈述是否正确、是否可迁移做一次语义准入检查。
本节据此说明旧实现实际做了什么，并未重新执行它，也未给旧轨迹补上今天设计的 Gate。
\end{sloppypar}

为什么不能简单补一行“失败就不写”？考虑一条失败轨迹：任务要求闭区间，程序却排除了右端点。
如果从中提取的是“检查端点是否包含”，失败来源仍可能产生有用提醒；如果提取的是“所有循环上界都加一”，同一个来源就可能产生错误泛化。
反过来，一次成功也可能只碰巧覆盖了特殊输入。从“本题通过”推不出“由它概括出的规则对后续题有效”。
这两个例子是机制说明，不是新增实验结果。

因此，需要分开保存来源运行的判定、候选经验的检查结果以及发布决定。
第一项回答当时做得怎样，第二项回答当前证据支持这条经验到什么程度，第三项回答允许它以什么范围影响后续运行。
一条经验可以被保存用于诊断却不参与检索；也可以只在注明条件的任务中试用。
是否采用这些状态要由具体协议决定，不能因为用了“经验记忆”这个名称，就认为检查与发布设施已经存在。

\section{版本、检查点、候选档案与血缘}
\label{sec:ch02-lineage}

Gate 决定之后，还剩一件事：把这一连串决定记下来。四个概念必须分开存，因为它们回答的是四个不同的问题。

\begin{table}[htbp]
\centering
\caption{版本、检查点、候选档案与血缘：存什么，为什么不能合并}
\label{tab:lineage}
\begin{tabular}{>{\raggedright\arraybackslash}p{2.2cm}>{\raggedright\arraybackslash}p{4.0cm}>{\raggedright\arraybackslash}p{4.0cm}>{\raggedright\arraybackslash}p{3.4cm}}
\toprule
概念 & 存什么 & 回答什么问题 & 合并了会丢什么 \\
\midrule
版本 $A_k$ & 初始系统 $A_0$ 或后续已采纳的版本 & 这个编号对应哪组组件 & — \\
检查点 & 可回到该状态所需的全部落盘 & 出事时能不能退回去 & 只有版本号没有检查点，回滚是空话 \\
候选档案 & 所有候选，\textbf{含被 reject 的} & 试过什么、为什么没被接受 & 失败分支消失，负面证据无法复核 \\
血缘 & 父版本、候选与组件变更的关联 & 某项修改从哪一版引入 & 多次更新后难以追溯修改 \\
\bottomrule
\end{tabular}
\end{table}

被 reject 的候选必须留档。这不是宽容，而是证据纪律：
“试过并且失败了”是判断候选表现的一类直接证据。它未必足以排除整条路线，但删掉它会妨碍故障复核，并增加重复无效尝试的风险。
开篇演示的运行目录提供了一个例子。每份 \path{trace.json} 保存一次回合的轨迹，\path{summary.json} 汇总各回合的标识与判定，并记录当时的 Git 提交号。该提交包含静态 Agent，却没有演示驱动脚本。六条轨迹仍可逐次核对验收和读写动作，要还原当时运行的完整程序，还缺驱动脚本在运行时的源码快照。

\section{从候选到运行版本：把发布做成可检查的状态转换}
\label{sec:ch02-publication}

一个候选通过了测试，运行进程却继续读取旧文件；另一个候选尚未通过测试，后台写入已经被检索器读到。
两种情况都可能发生在``生成、验证、发布''三步看似齐全的系统中。
缺少的是这三步之间的状态约束：验证究竟检查了哪个对象，发布指向哪个对象，执行器又实际加载了哪个对象。
本节给出本书建议的工程设计，不声称现有实验台已经实现了全部发布设施。

\subsection{候选内容与候选状态分开保存}

候选的内容可以是一条经验、一个技能文件或一组工作流节点；候选的状态则说明它处于提出、检查、拒绝、观察还是已发布阶段。
把两者放在同一个可随意覆盖的文件里，会制造一种难以复核的情况：昨天通过检查的是旧内容，今天文件名没变，内容却已经被再次修改。
一张写着``通过''的记录无法再对应到当前对象。

更容易检查的做法是，内容一旦送检便固定下来，后续修改产生新的候选标识。
状态记录引用固定内容的校验信息，同时保存使用的验证协议版本。
这样，同一个名字下的第二次修改不会借用第一次的通过状态。
哈希在这里承担的是对象绑定，不承担质量认证；一个错误程序也有稳定哈希。

为了说明约束，可以把候选表示为
\begin{equation}
c=(\mathrm{id},\ \mathrm{parent},\ \mathrm{artifact},\ \mathrm{evidence},\ \mathrm{status}).
\label{eq:ch02-candidate}
\end{equation}
父版本说明它从哪里来，工件指向实际内容，证据指向检查记录，状态决定它现在能否被执行器选用。
这只是本书的设计记号，不要求读者采用某种特定存储格式。
关键约束是：更换工件或验证协议以后，旧证据不能无条件为新对象背书。

\subsection{验证结束到发布之间，父版本可能已经变化}

设修改器从版本甲出发产生候选乙，检查期间另一个候选丙已经被发布。
乙的测试比较的是``甲与乙''，现在若直接让乙替换丙，就可能丢掉丙带来的修复。
更微妙的情况是，乙只保存了局部差异，系统把它套到丙上；组合后的对象既不是甲，也不是检查过的乙。
此时所有检查都可能真实执行过，真正上线的版本却没有被检查。

最小防线是在发布时再次核对父版本。
当前运行指针仍指向甲，乙才可以按原证据进入下一步；如果已经改变，就暂停这次发布并重新构造待检验组合。
重新检查不必机械地重跑全部任务，但缩小检查范围必须有依赖分析支持，不能只因改动看起来很小而跳过。
对小型实验，串行更新往往比处理多个并发候选更容易解释。

这里需要区分两种``串行''。
候选生成可以并行，节省等待时间；正式发布仍可以保持单一顺序，保证每次接受都有明确父版本。
如果连候选生成也依赖其他候选的中间结果，血缘就应记录这些额外依赖，而不能只留一个看起来整齐的父编号。
血缘记录的目的是描述真实关系，不是为了把所有过程强行画成一条链。

\subsection{发布指针与加载确认不是同一件事}

把一个配置项改成新版本，只证明控制侧发出了切换请求。
运行进程可能缓存了旧检索库，某个工具服务可能尚未重载，正在执行的任务也可能仍沿用旧提示。
因此结果记录应包含执行器实际加载的组件标识，而不只包含调度器打算发布的编号。
否则，发布后的失败可能被错误归给一个根本没有被加载的新版本。

最简单的实验可以在每次新任务开始时读取一个完整的组件清单，然后把清单绑定到本次轨迹。
任务过程中如果允许切换组件，则每次切换都应留下事件；如果不允许，就应固定到该任务结束。
两种策略都可以研究，但不能让同一条轨迹中间悄悄换了模型或技能，却仍用单一版本号解释全部动作。
复核者需要知道每一步由哪组条件产生，尤其是出现失败或成本突增的步骤。

加载确认也不能只看文件名。
同名文件可能被替换，远程索引可能在本地配置未变时刷新，模型别名也可能指向不同后端。
能固定的内容保存校验信息，不能固定的服务状态明确记录身份和时间限制。
如果某个外部条件无法重建，报告应缩小复现主张，而不是用一个更长的版本字符串掩盖缺失。

\subsection{影子运行必须隔离外部效果}

候选输出没有展示给用户，并不意味着影子运行没有影响外部世界。
一个候选可能调用发送消息、写数据库或提交订单的工具；即使最终答案被丢弃，工具效果仍然已经发生。
因此影子运行的约束应施加在工具执行与外部写入上，而不只是施加在最终回答上。
哪些工具可以只读重放，哪些必须模拟，哪些无法安全并行，都属于部署协议的一部分。

在编程题上，隔离执行器可以为正式版本与影子版本分别准备任务副本，让各自的文件写入留在自己的目录里。
对于依赖实时库存、共享配额或外部交易的任务，候选和正式版本可能无法面对完全相同的环境状态。
这时应记录比较条件的差异，避免把环境变化误判为版本效果。
影子运行增加了对真实任务分布的观察机会，但不会自动提供一个无偏且无副作用的对照实验。

灰度发布还引入了路由规则。
如果困难任务总被送给旧版本，简单任务总被送给新版本，两个版本的平均表现便不能直接相减。
随机路由、按任务类型分层或配对重放各有适用条件，选择哪一种必须在比较之前明确。
本章只规定运行记录要保留路由依据；具体统计设计由第\ref{ch03}章处理。

\subsection{恢复检查要回到可执行状态}

一个旧文件仍在磁盘上，不等于系统一定能恢复到旧状态。
旧技能可能依赖已经升级的工具接口，旧记忆索引可能与当前编码方式不兼容，旧参数也可能需要不同的加载配置。
因此检查点应保存恢复所需的依赖关系，并通过一次受控加载确认它仍然可执行。
仅记录版本号，可以帮助寻找对象；无法替代恢复演练。

恢复的对象也需要分层。
本地提示和代码可以切回，外部服务上已经发生的写入则可能需要补偿操作；已经给出的建议无法从使用者记忆中撤回。
在实验报告里，恢复成功应说明恢复了什么：组件文件、下一任务的执行条件，还是外部状态也已复原。
不说明对象的``一键回滚成功''，很容易把局部恢复说成完整恢复。

为减少这种含混，可以在发布前写一条恢复检查：从哪个检查点启动，用哪个只读任务检查基本功能，哪些外部效果仍未处理。
这个检查不必每次都很庞大，但应能够实际运行。
如果恢复本身也依赖模型临时生成指令，恢复过程就引入了新的不确定性，需要另外留痕与验证。
关键路径上的旧版本、依赖清单和恢复入口，不应只存在于修改器的临时上下文里。

\section{预算也是状态：配额到底属于谁}
\label{sec:ch02-budget-state}

预算常被写成运行函数的一个参数，看起来只是数字。
一旦系统有跨任务记忆、技能修订或候选搜索，预算实际还包含生命周期：它何时创建、被谁消耗、在哪个边界重置。
这条生命周期如果没有写进协议，同一个数值就可能代表完全不同的资源限制。

设协议允许一条技能序列拥有一次额外修订机会。
实现若在每次失败时重新创建计数器，那么每个失败点都能获得一次修订；实现若在整个序列开始时创建并一直携带计数器，所有失败点才真正共享这一次机会。
两种实现都可以作为设计选择，区别在于它们回答不同问题，不能使用同一个预算名称而不交代边界。
第\ref{ch08}章的 E05 修订正是把边界从文字约定落实到序列对象。

预算记录至少要区分授权上限、预留额度、实际消耗与未知成本。
某次模型调用已发出但网络连接中断，客户端可能没有得到 token 账；此时不能将该次成本记成零。
同样，提前失败只消耗较少调用，并不意味着整个方法更高效，因为它可能没有完成原定任务。
预算与完成条件应一起报告，才能解释少花资源究竟来自节省还是未执行。

还有一类易漏成本发生在候选之间。
系统生成多个修改，只保留最后成功的那个；若只给赢家记账，失败候选与筛选反馈便消失了。
对比较整个更新流程的问题，这些成本都属于流程；对研究单次部署推理的问题，可以另列为离线成本，但必须保留入口。
把离线与在线分开是口径说明，不是从总账中删除其中一部分。

恢复与重试也会改变预算解释。
在同一序列的预算边界内，重启不能重置可用配额。应从持久账恢复剩余额度，并保留已发生的消耗与尚未结清的预留。仅把旧消耗补入副账，仍可能让新计数器再次花出完整额度。若原账无法恢复，就无法完整核对总成本，也无法确认原配额仍得到遵守。
若不能确认外部请求是否已执行，直接重发还可能造成重复效果。
这就是为什么运行编号、序列编号、请求记录与版本编号应彼此关联，却不能互相替代。
版本回答用了哪套系统，运行回答做过哪次尝试，预算回答为这些尝试付出了什么。

一个可复核的更新闭环，因此不只包含成功发布。
它还包含提出后被拒绝的候选、验证时发现的冲突、未完成的外部调用，以及恢复后仍未解决的效果。
这些记录看起来不如成功率曲线整齐，却是解释系统是否持续改进所必需的边界条件。
在它们缺失时，下一步通常应该先补记录与状态约束，再增加修改器的复杂度。


\section{沿一条更新记录定位故障}
\label{sec:ch02-walkthrough}

前面的对象与边界可以放进同一个例子里检查。下面是本书构造的设计示例，不是一次新增模型实验，也不提供成功率证据。任务是维护一个区间查询工具：调用者给出若干起止位置，工具返回相应片段。执行器可以修改任务目录里的实现文件；跨任务可更新的对象只有技能库 $L$。验收规则、任务路由和模型参数都不允许由修改器改写。例子的重点是一次故障应该沿哪条记录追查，而不是推荐一种通用的区间算法。

\subsection{先写清楚这次任务与下次任务的边界}

第一项任务约定右端点不包含在结果内。执行器完成了程序，来源任务通过检查，修改器据此提出一个技能候选。这个候选可以包含函数、适用条件以及来源记录，但不能把来源任务的全部验收文件打包进去。否则所谓复用可能只是下一次再次拿到同一份题目与答案，无法区分技能与题目残留。

当前任务通过，回答的是程序在当前验收条件下能否完成这项任务；候选能否入库，还取决于它满足哪一种准入政策。最窄的政策只要求来源任务通过、工件完整且不包含受保护材料。这种政策允许我们研究后续迁移是否成功，但不能把入库状态解释为已经证明迁移成功。更强的政策可以另外设置兼容性检查或开发任务，不过这些检查的成本和信息必须计入候选选择过程。

下一项任务重新建立工作目录。初始文件来自新任务，技能来自明确发布的库版本，两者不能通过同一个临时目录混在一起。加载记录因此至少要让复核者分辨：哪个文件是题目给出的，哪个文件由当前执行器生成，哪个文件来自上一任务留下的持久组件。三种来源都可能是合法输入，但合法性取决于协议，不能从文件现在位于哪里反推。

设下一项任务改为包含右端点。如果技能只有一段代码，没有注明端点约定，检索命中以后执行器就可能把旧约定带进新任务。此时文件复位可以完全正确，技能加载也可以完全正确，错误仍然发生。它属于持久物适用范围与读取方式的问题，不能归为环境污染。这个区别决定了修复方向：重新创建临时目录不会解决技能语义不匹配，修改技能也无法修复错误的目录复用。

\subsection{把首次失败保留下来，再决定是否修订}

假定协议允许执行器在一次任务内读取报错并重试。这种重试改变当前上下文和任务产物，并不自动产生新的技能版本。如果任务最后通过，只能说明执行器在规定的任务内预算下完成了任务；技能库是否被改过，要另查持久写入记录。不能因为最终答案已经修好，就在实验汇总里补上一条并未发生的技能改进。

如果协议还允许修改器修订库中技能，任务内修正与跨任务修订就形成两条不同的路径。前者可能只把本题调用方式改对，后者可能给技能增加明确的端点参数，或者收紧适用范围。收紧范围也可能是合理更新：一个不再被误用的窄技能，未必比覆盖更多任务却容易误用的技能差。评价这种更新时，需要同时观察正确性、被选用的机会与放弃复用后发生的额外成本。

首次复用失败的记录不能被修订后的成功覆盖。保留两次结果，才能回答两个问题：原技能直接拿来用会怎样；系统发现失败以后能否恢复。如果只保留后者，读者看到的便是已经享受额外预算的成绩，却可能误以为这是首次复用表现。第\ref{ch08}章据此区分首轮结果与修复结果，本章强调它们首先是两个事件，随后才是两个统计口径。

失败也不必总触发修订。任务可能超出技能的声明范围，验收器可能遇到基础设施错误，剩余额度也可能不足。选择器拒绝修订时，应记录拒绝的条件；缺少这个条件，后续复核只能看到一个没有新版本的失败，却无法知道是系统没有提出更新，还是更新被规则挡住了。两者对应的改进方向并不相同。

\subsection{一次技能变化如何经过四个角色}

执行器负责产生可检查的任务轨迹。对这个例子而言，至少要知道它读取了什么技能、调用了什么工具、形成了什么代码、实际运行过哪些检查。最终回答里声称“已测试”，不等于轨迹里存在测试事件。反馈器应从真正发生的执行中取得结果，并把基础设施失败与程序行为失败区分开来；如果二者暂时无法区分，也应保留这个不确定状态。

修改器接收的是一组证据，而不是一句笼统的“做错了”。报错位置可以帮助定位问题，但不能独立确定问题原因。例如结果少了末尾元素，既可能来自端点约定不匹配，也可能来自输入预处理错误。修改器据此提出修订，仍然只是一个待检验假说。把失败解释写得很流畅，不会自动提高该解释的证据等级。

选择器检查候选是否符合已经规定的接受条件。它可以拒绝一个看起来有道理却没有通过检查的修订，也可以要求补做检查。接受意味着该候选满足本次发布政策，不意味着它从此适用于所有任务。政策只检查了端点行为，证据范围就到端点行为；若声称同时提高了异常输入处理能力，还需要相应检查支持。

这四个角色可以由不同程序实现，也可以共享同一个模型。共享实现并不取消职责划分，但会改变独立性。执行器和修改器使用同一模型时，不能把两个角色的相互同意当成两份独立证据；反馈器若也由该模型生成，自洽的错误还可能在多个角色之间循环。本书的角色图描述信息流与权限，不保证各角色的错误彼此独立。

\subsection{从结果异常倒查，而不是先给机制贴标签}

现在假定发布后出现了一个失败。排查应先确定这条失败属于哪个运行实例，再查该实例实际加载了什么。若它仍在使用旧库，就先处理发布与加载不一致；若它加载了新库，再查新库是否与送检工件一致；若工件一致，继续检查本次任务是否落在已声明范围内。这个顺序能够避免在没有确认运行条件时，就开始讨论新技能是不是退化了。

工件一致而任务超出范围，首先说明适用判断需要检查。工件一致、任务也在范围内，却没有覆盖到当前失败条件，则暴露了检查覆盖的缺口。若检查已经包含这个条件，实际结果却不同，就要核对依赖、工具版本、输入转换和随机性。几种情形最后都表现为失败，但它们要求的修复不在同一层。

这条倒查链也适用于成功。新版本取得好成绩，若运行记录显示它并未加载新技能，该成绩就不能作为技能带来收益的证据；若它加载了技能，同时任务已经换成更简单的一批，也不能仅凭平均分差归因。成功结果值得保存，但保存不等于已经完成解释。归因需要结果、实际条件与对照设计相互对应。

诊断可以提出多个相容假说。例如加载无误且任务范围相符时，失败仍可能由技能缺陷或执行器错误调用引起。在进一步检查之前，把两者都列为待检验原因，比直接写成“技能遗忘”更准确。术语应当压缩已经确认的机制，不能替代对机制的确认。

\section{框架如何落实为一个可反驳的比较}
\label{sec:ch02-testable-comparison}

运行记录把过程串起来，对照设计才让改进主张有机会被推翻。仍用区间工具说明：若研究问题是“发布修订后的技能，是否改善后续任务表现”，需要比较的对象是旧技能版本与新技能版本，而不只是修订前后两个时间段。时间段之间可能同时发生任务、模型服务或工具条件变化；这些变化会使前后分数差混入其他原因。

\subsection{先冻结要比较的对象}

一个最小的开发比较，可以把同一批开发任务分别交给旧版与候选版，在各自干净的任务环境中执行。两者使用相同的基本模型配置、相同工具能力与预先说明的预算规则。这样仍不保证一切随机因素完全相同，但至少能够说明计划改变的是哪个组件，以及哪些条件被努力保持一致。这里说的是设计要求，不是本书已经完成了一次新的配对实验。

比较中若允许候选继续学习，问题就变了。旧版是固定对象，新版却按任务顺序持续改变，最后的分数比较的是固定系统与在线更新流程。这个问题同样有价值，不过需要固定更新规则，记录每个任务实际面对的版本，并说明任务顺序如何产生。不能用静态版本比较的名称，报告一条持续变化的学习轨迹。

开发比较还参与候选选择。修改器看过结果再改，第二个候选便利用了第一轮反馈；即使每轮都叫“测试”，这些数据也已经参与开发。候选选择结束后，若要报告独立评价，就需要另外冻结评价入口和信息使用规则。独立性来自真实的信息边界，不来自目录名称中是否出现了 test 字样。

\subsection{预算相等之前，先说明什么资源相等}

给两版相同的调用次数，可以控制调用机会，却不能保证 token、时间或费用相同。新技能可能缩短搜索，也可能因为注入内容更长而增加输入成本。若研究关心的是固定调用预算下的任务表现，调用数相等是有意义的条件；若研究关心总费用，则还需要另一套费用口径。一个实验不必同时解决所有问题，但应避免让一种相等冒充所有资源都相等。

修改本身也有成本。只比较发布后的任务执行，可以研究已经拥有新技能时的边际表现；把候选生成、筛选与失败尝试加回来，才更接近采用这一更新流程的总成本。两份账都可以报告。关键是不能在需要证明整个流程划算时，只引用已经摊掉构建成本后的结果。

对于无法准确计量的外部成本，协议应规定如何报告缺失。请求已发出但未收到响应，不能因为轨迹没有 token 字段就补零；有些耗时来自共享后端排队，也不能全部归为算法本身。缺失标记会让总账看起来不那么完整，却保住了账的含义。重新测量是否必要，取决于这些缺失是否足以改变所关心的比较。

\subsection{冻结的是评价职责，而不是禁止发现评价错误}

默认冻结验证器，不代表验收程序永远不会出错。假如区间任务的验收断言把端点约定写反，坚持使用错误断言并不能提高科学性。合理处理是保存旧验收版本及其影响，修正错误后建立新的评价版本，再按新版本重算受影响结果。修正前后的分数不能无说明地接成同一条曲线。

评价错误的发现者也需要留下来源。若修改器在看见失败后提出“测试错了”，这只是一个待核验主张；应该依据任务定义和独立检查确认，而不是让修改器自己批准更改判分规则。反过来，外部发现的明确验收错误也不应因为协议原先冻结而被忽略。冻结的作用是阻止待评价系统暗中移动标尺，不是把已知错误变成不可纠正的规则。

修正后可能出现两种范围。若错误只影响一个明确子集，可以重算该子集并说明未受影响部分；若错误改变了任务语义或候选选择过程，已有候选可能已经针对旧规则被筛选，单纯重新打分无法消除这段选择历史。此时需要承认旧实验回答的问题已经受限，必要时用新设计重新开始。新运行编号不能抹去旧过程，旧结果也不能为新设计预先提供结论。

\subsection{把结论写到能够承担的程度}

完成检查后，结论可以依次回答三件事：持久组件是否实际改变，改变后的对象是否被后续任务使用，以及在声明的条件下出现了怎样的结果。这三件事需要的证据不同。前两件主要依赖工件、加载与轨迹记录；第三件还依赖任务划分、比较方式、预算和统计不确定性。工程记录齐全，不自动使小样本收益变得可靠。

一个候选被拒绝，也可以产生有用的结论：它没有通过哪项检查，失败条件能否重现，现有证据排除了哪种解释。但如果失败原因仍不清楚，就只能报告候选在当前条件下失败，不能据此宣布整类方法不可行。候选级、流程级与方法类别级的主张，所需证据范围逐级扩大。

读者据此可以把一篇系统报告拆成两层阅读。先沿版本与运行记录确认它实际做了什么，再沿比较设计确认数据支持什么。缺少前一层，分数可能找不到对应对象；缺少后一层，完整轨迹也只能说明发生过一个过程。本章的框架让这两层能够连接起来，第\ref{ch03}章进一步规定如何估计差异，以及什么结果可以回答预先规定的比较问题。


\section{条件变化：换任务、换模型、换反馈器之后}
\label{sec:ch02-conditions}

一套框架有没有用，看它在条件变化之后还剩下什么。

\textbf{换任务分布。}
上一次留下的文件、缓存与配额可能帮助也可能干扰新任务；影响方向必须测量，不能预设为单向恶化。
留存判据与起点判据的定义不变，但仍须核对新任务实际继承了什么；复位检查必须\textbf{每个任务执行一次}——
演示里的条件二只在第一次复位，后续起点差异必须进入结果解释。

\textbf{换模型。}
待检验的记忆 $M$ 可以保持不变，但基础模型是 $A_k$ 中的 $\theta$，换模型仍然改变了系统版本。留存判据与起点判据的定义仍然适用，但新模型实际获得哪些持久信息、从什么状态出发，仍须重新核对。
“这个持久物还有没有用”则须另行检验——一条写给模型甲的经验，在模型乙身上可能反号。
第 \ref{ch05} 章的 E03c 让两个读取者分别使用同源经验，并在每个读取者内与其无记忆基线比较。要判断记忆的收益是否随读取者改变，还须估计这两个差值之差及其不确定性。

\textbf{换反馈器。}
反馈器改变时（例如用模型自评替换执行反馈），应先确认信号用于产生候选，还是也进入 Gate。若 Gate 所用的有效信息减少，可能出现两类错误：该发布的不发布，不该发布的发布了。具体错误率与任务和反馈实现有关，不能从反馈名称直接推出。
$\mathcal V$ 冻结的价值在这里变成可操作的底线——
至少还有一套不参与修改的验收在，它让“判错了”可以被事后发现，而不是永远沉在账里。
第4章比较不同反馈方案在历史E02设置中的表现。各反馈臂分别训练并构建经验，观察到的差异属于完整方案，不能直接等同于反馈强度的净效应。

\textbf{换可更新集合。}
这改变的是允许修改的组件边界，三层时间尺度与四个角色仍可继续使用。
每扩大一档，差异定位、三条判据与复位检查都要重做一遍——这就是公理 2 的代价，也是第 12、14 章各自的起点。

% ---------- 6. 选择原则 ----------
\begin{principlebox}
判断一次变化值不值得记到系统上，按三件事依次检查：
\textbf{它在任务结束后还在不在}（留存）、\textbf{下一个任务有没有因为它从不同地方出发}（起点）、
\textbf{这个不同能不能被一个不参与修改的机制检验}（独立检验）。
前两条加上明确的系统边界用于确认受控持久更新，第三条为判断效果提供独立证据。缺少效果证据时，应记录为已经发生但效果未明的更新，不能把它自动归为环境残留。
把可更新集合扩大一档，就等于把这三件事重新做一遍。
\end{principlebox}

\begin{pitfallbox}
\textbf{反例一：环境残留冒充学习。}
本章开头的演示里，两种条件的验收结果都是 \measured{3/3}，
但共用目录的后 \measured{2} 次运行都读到了第一次写下的修正。
本例未测得分数差；更难任务是否出现差异，需要另外运行，不能从起点污染直接推算。

\textbf{反例二：把 $\mathcal V$ 放进可更新集合。}
若任务与评分方式同时变化，式~\eqref{eq:delta} 的差值不再单独表示同一目标下的系统收益。
需要额外冻结共同评价条件，才能分开解释系统与标尺的变化。

\textbf{反例三：只留版本号，不留检查点与候选档案。}
出事时退不回去，失败分支也不复存在；下一次团队会重新试一遍同一条死路。
\end{pitfallbox}

\section*{习题}
\begin{enumerate}[label=2-\arabic*.,ref=2-\arabic*]
  \item\label{ex:ch02-1} 选一个你熟悉的 Agent 系统，按表~\ref{tab:roles} 指出它的修改器搜索空间 $\mathcal S$ 是什么。
        如果它的 $\mathcal S$ 里包含验收器，式~\eqref{eq:delta} 在它身上还成立吗？
  \item\label{ex:ch02-2} 本章的演示里两种条件的通过率相同。若要检验目录残留是否影响通过率，你会怎样预先选择任务、确定运行次数并安排对照？
        区分起点一致性检查与通过率差异检验，并说明无差异结果应如何报告。
  \item\label{ex:ch02-3} 三级 Gate 中，shadow 与 accept($r$) 的信息差别是什么？
        如果一个系统只有 reject 与 accept，它的哪一类决策做不了？
  \item\label{ex:ch02-4} 举一个你见过的“环境残留”实例（缓存、临时文件、数据库状态、外部服务配额都算），
        并说明你会用什么检查把它与持久更新分开。
\end{enumerate}

\section*{常见误解}
\begin{itemize}
  \item \textbf{“起点变了一点没关系，反正分数是靠系统拿到的”。}
        起点若偏离声明的比较条件，前后差值就可能同时包含残留状态的影响，不能仅凭该差值归因于受控系统更新。
  \item \textbf{“Gate 就是验收通过就发布”。}
        那是二分；它做不了“值得观察但不值得发布”这一类候选；这类候选是否常见，需要具体系统的记录支持。
  \item \textbf{“被 reject 的候选可以不记”。}
        失败分支是证据。删掉它，负面结果就无法复核，同一条路线会被重复尝试。
\end{itemize}

\section*{进一步阅读}
\begin{itemize}
  \item \cite{self-improvements-survey-2607}：状态分解与 operational scaffold 的出处，
        本书的 $A_k$ 沿用其组件化描述思路，并单列评价设置；读它可以比较两套对象边界。
  \item \cite{self-evolving-comprehensive-2508}：第3.5节给出修改器 $(\mathcal S,H)$ 分解与式(1)的搜索目标。
        其图3的四个组件是输入、Agent系统、环境与优化器；本章另按执行、反馈、修改与选择划分职责，两套划分不能逐项等同。
  \item \cite{harness-eval-2607}：观察四种改进路径怎样设置对照；第3章区分这些路径与本书按持久范围组织的四格。
  \item \cite{rsi-survey-2607}：把 bounded 与 open-ended 的划界讲清，
        本章主要处理可明确记录对象、预算与评价条件的更新；更开放的系统同样需要这些记录，但仅靠这些记录不足以定义其长期目标。
\end{itemize}
